\documentclass[]{article}

\usepackage{amsmath}
\usepackage{multirow}
\usepackage{natbib}
\usepackage{graphicx} 


\begin{document}

Our analysis of the Lagrangian trajectories of vortex structures reveals a robust superdiffusive transport regime. We characterize this behavior by examining the Mean Squared Displacement, the temporal correlations of vortex velocity, and the statistical distribution of their step sizes. We further investigate the physical origins of this transport by quantifying the coupling between vortex motion and the background fluid flow, as well as the anisotropy of the motion relative to the vortex's own orientation.

\subsection{Superdiffusive nature of vortex transport}
The primary diagnostic for the transport regime is the Mean Squared Displacement (MSD), which we computed for the ensemble of vortex trajectories. A power-law fit to the MSD, $\text{MSD}(\tau) \propto \tau^\alpha$, yields a diffusive exponent for the total displacement of $\alpha_{\text{total}} = 1.815 \pm 0.009$ ($R^2 = 0.999$). This value is significantly greater than the exponent for classical diffusion ($\alpha=1$) and ballistic motion ($\alpha=2$), placing the vortex transport firmly in the superdiffusive regime. This indicates that vortices are transported through the turbulent flow far more efficiently than would be expected from a simple random walk.

To investigate potential anisotropy in the transport, we analyzed the MSD for each Cartesian component separately. The analysis revealed distinct diffusive exponents for each direction:
\begin{itemize}
    \item $\alpha_x = 1.724 \pm 0.016$ ($R^2 = 0.996$)
    \item $\alpha_y = 1.871 \pm 0.003$ ($R^2 = 1.000$)
    \item $\alpha_z = 1.849 \pm 0.009$ ($R^2 = 0.999$)
\end{itemize}
While all components exhibit strong superdiffusion, there is a noticeable anisotropy, with the transport being most pronounced in the y-direction and least pronounced in the x-direction. This suggests that the large-scale structure of the turbulent flow may impose a preferential direction for vortex transport.

\subsection{Temporal correlations and the nature of the random walk}
To understand the mechanism driving the observed superdiffusion, we investigated the temporal characteristics of the vortex motion. We first computed the Velocity Autocorrelation Function (VACF), which quantifies the 'memory' in a vortex's velocity. The VACF was found to decay slowly, with a characteristic correlation time (the time to decay to $1/e$) of approximately $\tau_{1/e} \approx 0.19$ simulation time units. This slow decay signifies that the vortex velocities are persistent over time; a vortex's motion at one instant is strongly correlated with its motion in the recent past. Such long-lived temporal correlations are a hallmark of a correlated random walk and are a primary driver for superdiffusive behavior.

Next, we analyzed the probability distribution of the trajectory step sizes, $|\Delta \mathbf{r}|$, to distinguish between different superdiffusive models. Specifically, we tested the hypothesis of a Lévy flight, which is characterized by a heavy-tailed distribution of step lengths. Our analysis shows that the step-size distribution is nearly Gaussian, with a kurtosis of less than 0.5. A Kolmogorov-Smirnov test confirms that the distribution is not statistically distinguishable from a Gaussian distribution ($p > 0.05$). This result allows us to rule out Lévy flights as the underlying mechanism. The combination of a slowly decaying VACF and a Gaussian-like step-size distribution strongly supports the conclusion that vortex superdiffusion is best described as a correlated, or persistent, random walk.

\subsection{Coupling with the background flow and transport anisotropy}
To connect the statistical properties of the transport to the underlying fluid dynamics, we quantified the coupling between the vortex motion and the surrounding flow. We computed the Pearson correlation coefficient between the components of the vortex centroid velocity and the interpolated fluid velocity at the centroid's location. The results, summarized in Table \ref{tab:vel_corr}, reveal a strongly anisotropic coupling.

\begin{table}[h!]
\centering
\caption{Pearson correlation (r) between vortex centroid velocity and local fluid velocity components. The p-value indicates the statistical significance of the correlation.}
\label{tab:vel_corr}
\begin{tabular}{lcc}
\hline
\hline
Component & Pearson r & p-value \\
\hline
$v_x$ & 0.139 & $4 \times 10^{-17}$ \\
$v_y$ & 0.790 & $\approx 0$ \\
$v_z$ & 0.059 & $4 \times 10^{-4}$ \\
Speed ($|\mathbf{v}|$) & 0.336 & $\approx 0$ \\
\hline
\end{tabular}
\end{table}

The correlation is exceptionally strong in the y-component ($r=0.790$) but weak in the x and z components. This anisotropy in coupling directly corresponds to the anisotropy observed in the MSD exponents, suggesting that the enhanced superdiffusion in the y-direction is driven by a stronger advective coupling to the large-scale flow along that axis. However, the overall directional alignment between the vortex and fluid velocities is moderate, with a mean cosine of the angle between them of 0.317. This indicates that while vortices are significantly advected by the local flow, their motion is not entirely slaved to it, pointing to the role of intrinsic vortex dynamics.

We further explored the intrinsic dynamics by examining the anisotropy of vortex motion relative to its own orientation, defined by the local vorticity vector $\boldsymbol{\omega}$. We decomposed each displacement step into components parallel ($\Delta\mathbf{r}_\parallel$) and perpendicular ($\Delta\mathbf{r}_\perp$) to $\boldsymbol{\omega}$. The analysis shows that the mean displacement perpendicular to the vorticity axis is significantly larger than that parallel to it:
\begin{itemize}
    \item Mean $|\Delta\mathbf{r}_\parallel| = 0.001800$
    \item Mean $|\Delta\mathbf{r}_\perp| = 0.002131$
\end{itemize}
The ratio of parallel to perpendicular mean displacement is 0.845, indicating that vortices preferentially drift in the plane perpendicular to their own axis. This is physically consistent with the behavior of elongated vortex filaments, which can be advected sideways more easily than they can be translated along their own length, an action that would require substantial internal shearing.

\subsection{Robustness of the superdiffusive regime}
To ensure the validity of our findings, we performed a series of tests to check the robustness of the measured superdiffusive exponent. First, we varied the threshold for the Q-criterion used to identify vortices, testing values of $Q > \mu_Q + k\sigma_Q$ for $k \in \{2.5, 3.0, 4.0\}$. The diffusive exponent $\alpha$ remained stable across all thresholds, confirming that our results are not sensitive to the precise definition of a vortex.

Second, we tested for convergence with respect to trajectory length by calculating $\alpha$ for subsets of trajectories with minimum lengths ranging from 20 to 100 snapshots. The exponent was found to be stable across this range, indicating that the observed superdiffusion is not a transient, finite-time artifact. Furthermore, by comparing the diffusive exponent calculated from the first and second halves of the trajectories, we found no significant difference, confirming that the transport process is in a statistically stationary state.

Finally, to isolate the intrinsic dynamics from passive advection, we re-calculated the MSD after subtracting a large-scale advective velocity field. The resulting residual trajectories exhibited an even stronger superdiffusive character, with $\alpha_{\text{residual}} \approx 1.93$. This counter-intuitive result provides strong evidence that the superdiffusive behavior is an intrinsic property of the vortex dynamics and not merely a consequence of being passively transported by a superdiffusive background flow.

\end{document}
                